\lecture{Tecnologie per Materiali Polimerici in Fase Liquida}{02~-~07~-~2026}
%
%
\section{Introduzione e Classificazione}
\begin{itemize}
    \item \textbf{Normativa:} La terminologia e i principi della fabbricazione additiva (FA) sono standardizzati dalla norma ISO/ASTM 52900:2021.
    \item \textbf{Dominio dei Polimeri:} I polimeri sono i materiali più utilizzati in assoluto nella fabbricazione additiva.
    \item \textbf{Stato di partenza:} Le tecnologie per i polimeri si classificano in base allo stato del materiale di partenza: Polvere (es. SLS, MJF), Solido/Filamento (es. FDM) e \textbf{Liquido}.
    \item \textbf{Tecnologie a base liquida:} Si dividono in \textit{Fotopolimerizzazione in vasca} (SLA, DLP) e \textit{Deposizione di gocce} (Material Jetting / Polyjet). Generalmente, le parti prodotte hanno tolleranze paragonabili alle lavorazioni convenzionali, ma presentano rugosità che possono richiedere finitura manuale.
\end{itemize}

\section{Vat Photopolymerization: Stereolitografia (SLA)}
La Stereolitografia è stata la prima tecnologia commerciale, inventata da Charles Hull (3D Systems, 1987).
\begin{itemize}
    \item \textbf{Principio di funzionamento:} Un laser UV solidifica selettivamente una resina fotopolimerica liquida contenuta in una vasca. La piattaforma (elevatore) si abbassa (o si alza) ad ogni strato, costruendo il pezzo dal basso verso l'alto.
    \item \textbf{Materiali (Fotopolimeri):} Si utilizzano resine termoindurenti. Storicamente si usavano \textit{Acrilati} (molto reattivi ma soggetti a ritiro/deformazione) ed \textit{Epossidiche} (più precise e resistenti, ma fragili e sensibili all'umidità). Le resine moderne sono ibride per bilanciare questi effetti.
    \item \textbf{Il problema del Ritiro (Curling):} Durante la fotopolimerizzazione, il volume si riduce. Questo restringimento tira gli strati sottostanti, accumulando tensioni residue che portano all'arricciamento dei bordi (deformazione).
    \item \textbf{Strutture di Supporto (Importante per l'esame):} Sono fondamentali per ancorare il pezzo e contrastare la deformazione. Nella SLA i supporti sono creati con lo \textit{stesso identico materiale} del pezzo e devono essere rimossi meccanicamente.
    \item \textbf{Vantaggi:} Eccellente precisione ($<0,1$ mm), strati sottilissimi ($0,05 - 0,025$ mm), ottima finitura superficiale, possibilità di usare resine trasparenti o per microfusione (calcinabili).
    \item \textbf{Svantaggi:} Necessità di post-trattamento (forno UV per completare la polimerizzazione), obbligo di lavorare a "vasca piena" (costo e gestione del liquido), tensioni residue.
\end{itemize}

\section{Vat Photopolymerization: DLP e CLIP}
\begin{itemize}
    \item \textbf{Digital Light Processing (DLP):} Simile alla SLA, ma al posto del laser utilizza un \textbf{proiettore digitale UV}.
          \begin{itemize}
              \item \textit{Vantaggio chiave:} Polimerizza l'intera sezione del piano in un solo istante (flash), rendendo il processo generalmente più rapido per strato rispetto al laser che deve tracciare il profilo.
          \end{itemize}
    \item \textbf{Continuous Liquid Interface Production (CLIP - Carbon 3D):} È un'evoluzione rivoluzionaria.
          \begin{itemize}
              \item Utilizza una finestra permeabile all'ossigeno sul fondo della vasca.
              \item L'ossigeno inibisce la polimerizzazione, creando una "dead zone" (Oxygen Inhibition Layer) liquida tra la finestra e il pezzo in formazione.
              \item Questo impedisce al pezzo di attaccarsi al fondo, permettendo un'elevazione \textbf{continua} e senza interruzioni (niente step di distacco per ogni strato), aumentando drasticamente la velocità di stampa.
          \end{itemize}
\end{itemize}

\section{Material Jetting: PolyJet e Multi Jet Modelling (MJM)}
A differenza delle tecnologie in vasca, qui il liquido viene depositato solo dove serve.
\begin{itemize}
    \item \textbf{Principio di funzionamento (PolyJet - Stratasys):} Testine di stampa (simili a quelle a getto d'inchiostro) depositano microscopiche gocce di fotopolimero acrilico, che vengono \textbf{immediatamente} solidificate da lampade UV poste sul carrello.
    \item \textbf{Strutture di supporto (Importante per l'esame):} I supporti vengono stampati con un \textit{materiale differente} (un materiale gel-like o cera) che è \textbf{solubile} (rimovibile con getto d'acqua o per fusione della cera nel caso del MJM di 3D Systems).
    \item \textbf{Materiali Digitali e Multi-materiale:} Il vero punto di forza del PolyJet. La macchina può depositare \textit{simultaneamente} materiali diversi (es. 6 materiali + 1 supporto). Miscelando gocce di materiali rigidi e flessibili (Rubber-like) "on-the-fly", si creano materiali compositi con proprietà meccaniche specifiche (es. precisi valori di durezza Shore) e colori diversi.
\end{itemize}
\begin{description}
    \item [Vantaggi] Nessuna vasca di resina liquida, nessun post-trattamento in forno necessario (i pezzi escono già polimerizzati), finitura liscia, supporti solubili (permette geometrie complesse interne), stampa multi-materiale/multi-colore per prototipi iper-realistici.
    \item [Svantaggi] Si usano materiali termoindurenti (quindi non hanno le stesse identiche prestazioni meccaniche dei termoplastici reali industriali), e il volume di lavoro non può essere sfruttato interamente impilando pezzi uno sull'altro (essendo una deposizione discreta).
\end{description}